% !TEX root = ../project_report.tex
\section{Discussion}

Recent rainfall substantially improves the short-term probability estimates, especially at lower thresholds. Calendar features describe recurring seasonal and daily patterns, while recent observations capture local rainfall persistence. Specifically, RF has the lowest observed threshold-average Brier among the classifier pipelines, with a small positive conditional advantage over GBT in A. GBT ranks 1~mm events competitively, and in B its Brier is lower at the two lowest thresholds. Historical frequency is strongest for calendar-only mean Brier. It is also worth noting that limitations like average risk agreement, coarse-bin calibration and ranking performance still exist.

The fitting-fraction comparison shows small mean changes at the chosen capacities and seed. RF still estimates within-bin distributions from the full training population, and LR/GBT expand each fitting anchor into seven binary examples. The results therefore do not quantify end-to-end savings from using less data. The depth-10 validation improvement is more promising than increasing the reference forest's fitting fraction, but it requires its own calibrated evaluation. Later reuse of 2024 limits confirmatory model-selection claims. Rolling chronological comparisons within the development years and evaluation in a subsequent untouched period would assess stability.

Coverage defines the population described. Task 1 retains 85 of 96 year-months; A retains 89.9\% of B's eligible 2024 examples. Its scores therefore apply to times with complete published history and future labels. At 1~mm, the 15\% A RF improves Brier over its same-setting historical-frequency baseline in every 2024 month (9.8--33.9\%) and at all 70 evaluated stations (22.1--41.9\%). These descriptive subgroup comparisons have no subgroup confidence intervals. Broader assessment should cover other thresholds and years. Neither the eight-year observation record nor a single evaluation year establishes long-term climate or model stability.

We checked the station-catalogue concern directly. Although offline B features take coordinates from observations and the query command uses the frozen training catalogue, all saved catalogues agree and every active 2022--2024 station-month matches them within numerical precision, with no within-month coordinate variation. Thus this policy difference does not change coordinates in the evaluated years. Future relocations require consistent catalogue updates. A live versioned archive would additionally establish which record version was available at each issue time.

\section{Conclusion}
The records show seasonal, diurnal and local rainfall variation across Singapore. Coverage checks and common-station comparisons support the descriptive findings, although changing station coverage and source gaps remain limitations. For Task 2, recent-rainfall RF is a useful evaluated reference, with gains concentrated at lower thresholds and limited heavy-rain probability skill. Historical frequency remains strongest for calendar risk at the evaluated thresholds. The close model and fitting-fraction comparisons support scoped conclusions about the retained pipelines; stronger validation and fine-resolution tail diagnostics are the next steps.

\par\noindent\begin{minipage}{\linewidth}
\section*{AI usage declaration}
OpenAI Codex assisted with code development, cloud execution, and report drafting.
\end{minipage}
